\subsection{The operationalisation of return, and a worked illustration}
\label{p11:sec:ev-operational}

The assessment of justice requires that the abstract notion of the return of value be operationalised into a quantity that could, in principle, be estimated; this subsection sets out the operationalisation and illustrates it on a deliberately simple example, under the operational concession throughout.

The operationalisation proceeds in four steps. First, the contributions of each party over a period are identified: the discrete acts of generation, the pieces of sustaining labour, the cared-for occasions, that each party contributes to the shared life. Second, to each contribution is assigned a value, in a common unit, proxying the value it generates for the shared life; the assignment is the most contestable step, and the concession applies most stringently here, since the value of an act of care is exactly the unsymbolised quantity the proxy cannot reach, and the assignment is an operational stipulation and not a measurement. Third, for each party the return is identified: of the value generated in the shared life, the portion that returns to that party, whether directly or through the relation's circulation. Fourth, the return to each party is normalised by the coupling-based baseline of \xref{p11:sec:ev-return}, the return the coupling structure would yield under unbiased propagation, to give the normalised return on which the assessment of skewness is conducted.

A worked illustration, with stipulated numbers, shows the quantities in operation; the numbers are illustrative and carry no empirical claim. Consider a deeply coupled relation of two parties, $A$ and $B$, over a period in which the shared life generates value, in the common unit, of 100. Suppose $B$ bears the larger share of the sustaining labour and generates 60 of this value, while $A$ generates 40. In a just cycle, given the depth of the coupling, the return to each would be approximately proportional to generation: $B$ would receive a return near 60 and $A$ a return near 40, the normalised returns of both being near 1, and the distribution of normalised return approximately symmetric. Suppose instead that the relation returns value disproportionately to $A$: of the 100 generated, $A$ receives a return of 65 and $B$ a return of 35. The normalised return to $A$ is then $65/40 \approx 1.6$ and to $B$ is $35/60 \approx 0.58$. The distribution of normalised return is now skewed: $A$ sits well above the baseline of 1, and $B$ well below it, in the low-return tail. The skewness of the normalised-return distribution, negative for $B$, is the statistical fingerprint of \xref{p11:sec:ev-skew}, and it is legible here despite the relation's possibly high mean satisfaction, since $A$, the over-returned party, may report great contentment, and $B$, habituated to bearing the sustaining labour for little return, may report a satisfaction only moderately depressed, so that the mean of the two is unremarkable while the skew is severe. The illustration shows the quantities the assessment computes; it does not show that they can be estimated easily in an actual relation, which, as \xref{p11:sec:ev-model} stated, remains a substantial empirical problem the paper does not solve.
